% Part: lambda-calculus
% Chapter: introduction
% Section: free-variables

\documentclass[../../../include/open-logic-section]{subfiles}

\begin{document}

\olfileid{lam}{syn}{fv}
\olsection{ازاد متغيرونه}

لامبډا حساب د تابعو په اړه دے، او تابعې د لامبډا تجريد له لارې
رامنځته کېږي۔ په شهودي ډول، $\lambd[x][M]$ هغه تابع ده چې قيمتونه
ئې~$M$ ټاکي، کله چې د تابعې آرګومېنټ~$x$ ته ورکړل شي۔ خو په~$M$
کښې د~$x$ هر مورد اړوند نۀ دے: که په~$M$ کښې بل تجريد
$\lambd[x][N]$ وي، نو په~$N$ کښې د $x$ موردونه د
$\lambd[x][N]$ دپاره اړوند دي، خو د $\lambd[x][M]$ دپاره نۀ دي۔
نو په $\lambd[x][M]$ کښې دننه يو لامبډا تجريد $\lambd[x]$ د~$M$
هغه د $x$ موردونه \emph{تړي} چې لا د بل لامبډا تجريد له خوا تړل
شوي نۀ وي---يعنې په~$M$ کښې د~$x$ \emph{ازاد} موردونه۔

\begin{defn}[ساحه]
که $\lambd[x][M]$ په يوه ترم~$N$ کښې واقع شي، نو د~$M$ اړوند
مورد د~$\lambd[x]$ د اړوند مورد \emph{ساحه} ده۔
\end{defn}
\textbf{د ژباړې د نسخې سمون (OLLAM-010):} سرچينې د لامبډا
تجريد د ساحې په تعريف کښې د چاپېره ترم N اړوند مورد راوړے و؛ ساحه
د تجريد د بدن M اړوند مورد دے، لکه د لاندې بېلګو او د متغيرونو د
تړلو له تعريفه چې څرګندېږي۔

\begin{defn}[ازاد او تړلے مورد]
  په يوه ترم~$M$ کښې د متغير $x$ مورد هله \emph{ازاد} دے چې د
  $\lambd[x]$ په ساحه کښې نۀ وي، او که وي نو \emph{تړلے} دے۔ په
  $\lambd[x][M]$ کښې د متغير~$x$ يو مورد هله او يوازې هله د لومړي
  $\lambd[x]$ له خوا تړلے دے چې په~$M$ کښې د $x$ همدا مورد ازاد وي۔
\end{defn}

\begin{ex}
  په $\lambd[x][x y]$ کښې $x$ او $y$ دواړه د $\lambd[x]$ په
  ساحه کښې دي، نو $x$ د $\lambd[x]$ له خوا تړلے دے۔ څرنګه چې $y$
  د هېڅ $\lambd[y]$ په ساحه کښې نۀ دے، ازاد دے۔ په
  $\lambd[x][x x]$ کښې د $x$ دواړه موردونه د~$\lambd[x]$ له خوا
  تړلي دي، ځکه دواړه په $xx$ کښې ازاد دي۔ په $((\lambd[x][xx])x)$
  کښې د~$x$ وروستے مورد ازاد دے، ځکه د $\lambd[x]$ په ساحه کښې نۀ
  دے۔ په $\lambd[x][(\lambd[x][x])x]$ کښې د لومړي $\lambd[x]$
  ساحه $(\lambd[x][x])x$ ده او د دويم $\lambd[x]$ ساحه د~$x$ له
  پای څخه دويم مورد دے۔ په $(\lambd[x][x])x$ کښې د $x$ وروستے
  مورد ازاد او له پای څخه دويم مورد تړلے دے۔ نو په
  $\lambd[x][(\lambd[x][x])x]$ کښې د $x$ له پای څخه دويم مورد د
  دويم $\lambd[x]$ له خوا او وروستے مورد د لومړي $\lambd[x]$
  له خوا تړلے دے۔
\end{ex}

په يوه ترم $P$ کښې د متغيرونو ټول موردونه کتلے شو او د ازادو
متغيرونو سټ ترې موندلے شو۔ دا سټ په $\FV{P}$ ښيو؛ طبيعي تعريف
ئې داسې دے:

\begin{defn}[د ترم ازاد متغيرونه] \ollabel{def:fv}
  د ترم د \emph{ازادو متغيرونو} سټ په استقرايي ډول داسې تعريفېږي:
  \begin{enumerate}
    \item $\FV{x} = \{x\}$ \ollabel{def:fv1}
    \item $\FV{\lambd[x][N]} = \FV{N} \setminus \{x\}$    \ollabel{def:fv2}
    \item $\FV{PQ} = \FV{P} \cup \FV{Q}$ \ollabel{def:fv3}
  \end{enumerate}
\end{defn}

\begin{prob}
  \begin{enumerate}
  \item په دې ترم کښې د $\lambd[g]$ او دواړو $\lambd[x]$ ساحې
    وټاکئ: $\lambd[g][(\lambd[x][g (x x)]) \lambd[x][g (x x)]]$۔
  \item په $\lambd[g][(\lambd[x][g (x x)]) \lambd[x][g (x x)]]$
    کښې ايا د متغيرونو ټول موردونه تړلي دي؟ هر يو د کوم تجريد له
    خوا تړلے دے؟
    \item $\FV{\lambd[x][(\lambd[y][(\lambd[z][x y]) z]) y]}$
    ومومئ۔
  \end{enumerate}
\end{prob}

\begin{explain}
ازاد متغير له باندنۍ نړۍ (\emph{چاپېريال}) سره د يوې حوالې په شان
دے، او هغه ترم چې ازاد متغيرونه لري د نيمګړي ټاکل شوي ترم په توګه
کتلے شو، ځکه چلند ئې د چاپېريال په ټاکلو پورې تړلے دے۔ مثلاً، په
ترم $\lambd[x][f x]$ کښې، چې آرګومېنټ~$x$ اخلي او د هغه آرګومېنټ
$f$ ورکوي، متغير~$f$ ازاد دے۔ د دې ترم قيمت په خپل چاپېريال پورې
تړلے دے، په تېره بيا په هغه چاپېريال کښې د $f$ په قيمت پورې۔

که پر دې ترم تجريد وکړو، $\lambd[f][\lambd[x][f
    x]]$ ترلاسه کوو۔ دا ترم نور د چاپېريال پر متغير $f$ تړلے نۀ
دے، ځکه اوس داسې تابع ښيي چې دوه آرګومېنټونه اخلي او د لومړي پر
دويم د تطبيق پايله ورکوي۔ په چاپېريال کښې د $f$ بدلون د دې ترم پر
چلند هېڅ اغېز نۀ لري؛ ترم يوازې هغه څۀ کاروي چې د آرګومېنټ په توګه
ورکول کېږي، نۀ په چاپېريال کښې د $f$ قيمت۔
\end{explain}

\begin{defn}[تړلے ترم، ترکيب کوونکے]
  هغه ترم چې ازاد متغيرونه ونۀ لري، \emph{تړلے ترم} يا
  \emph{ترکيب کوونکے} بلل کېږي۔
\end{defn}

\begin{lem}\ollabel{lem:fv}
  \begin{enumerate}
    \item \ollabel{lem:fv-abs} که $y \neq x$ وي، نو $y \in
      \FV{\lambd[x][N]}$ هله او يوازې هله چې $y \in \FV{N}$ وي۔
    \item \ollabel{lem:fv-app} $y \in \FV{PQ}$ هله او يوازې هله
      چې $y \in \FV{P}$ يا $y \in \FV{Q}$ وي۔
    \end{enumerate}
\end{lem}

\begin{proof}
  تمرين۔
\end{proof}

\begin{prob}
  \olref[lam][syn][fv]{lem:fv} ثابت کړئ۔
\end{prob}

\end{document}
